\documentclass[10pt,a4paper]{amsart}
\usepackage{amsmath,amssymb,amsthm,graphicx}
% For commutative diagrams we use Xy-pic
\usepackage[all]{xy}

% We use 2cell for 2-commutative diagrams.
\xyoption{2cell}
% Theorem environments.
%
\theoremstyle{plain}
\newtheorem{theorem}[subsection]{Theorem}
\newtheorem{proposition}[subsection]{Proposition}
\newtheorem{lemma}[subsection]{Lemma}

\theoremstyle{definition}
\newtheorem{definition}[subsection]{Definition}
\newtheorem{example}[subsection]{Example}
\newtheorem{exercise}[subsection]{Exercise}
\newtheorem{situation}[subsection]{Situation}

\theoremstyle{remark}
\newtheorem{remark}[subsection]{Remark}
\newtheorem{remarks}[subsection]{Remarks}

\numberwithin{equation}{subsection}

% Macros
%
\def\lim{\mathop{\mathrm{lim}}\nolimits}
\def\colim{\mathop{\mathrm{colim}}\nolimits}
\def\Spec{\mathop{\mathrm{Spec}}}
\def\Hom{\mathop{\mathrm{Hom}}\nolimits}
\def\Ext{\mathop{\mathrm{Ext}}\nolimits}
\def\SheafHom{\mathop{\mathcal{H}\!\mathit{om}}\nolimits}
\def\SheafExt{\mathop{\mathcal{E}\!\mathit{xt}}\nolimits}
\def\Sch{\mathit{Sch}}
\def\Mor{\mathop{\mathrm{Mor}}\nolimits}
\def\Ob{\mathop{\mathrm{Ob}}\nolimits}
\def\Sh{\mathop{\mathit{Sh}}\nolimits}
\def\NL{\mathop{N\!L}\nolimits}
\def\CH{\mathop{\mathrm{CH}}\nolimits}
\def\proetale{{pro\text{-}\acute{e}tale}}
\def\etale{{\acute{e}tale}}
\def\QCoh{\mathit{QCoh}}
\def\Ker{\mathop{\mathrm{Ker}}}
\def\Im{\mathop{\mathrm{Im}}}
\def\Coker{\mathop{\mathrm{Coker}}}
\def\Coim{\mathop{\mathrm{Coim}}}

% Boxtimes
%
\DeclareMathSymbol{\boxtimes}{\mathbin}{AMSa}{"02}

%
% Macros for moduli stacks/spaces
%
\def\QCohstack{\mathcal{QC}\!\mathit{oh}}
\def\Cohstack{\mathcal{C}\!\mathit{oh}}
\def\Spacesstack{\mathcal{S}\!\mathit{paces}}
\def\Quotfunctor{\mathrm{Quot}}
\def\Hilbfunctor{\mathrm{Hilb}}
\def\Curvesstack{\mathcal{C}\!\mathit{urves}}
\def\Polarizedstack{\mathcal{P}\!\mathit{olarized}}
\def\Complexesstack{\mathcal{C}\!\mathit{omplexes}}
% \Pic is the operator that assigns to X its picard group, usage \Pic(X)
% \Picardstack_{X/B} denotes the Picard stack of X over B
% \Picardfunctor_{X/B} denotes the Picard functor of X over B
\def\Pic{\mathop{\mathrm{Pic}}\nolimits}
\def\Picardstack{\mathcal{P}\!\mathit{ic}}
\def\Picardfunctor{\mathrm{Pic}}
\def\Deformationcategory{\mathcal{D}\!\mathit{ef}}
\usepackage{cleveref}
\crefname{originallemma}{original lemma}{original lemmas}
\Crefname{originallemma}{Original lemma}{Original lemmas}
\makeatletter
\newlabel{lemma-integral-B}{{04DR}{source}{}{sourcepointer-1}{}}
\newlabel{lemma-integral-B@cref}{{[originallemma][0][]04DR}{[1][0][]source}{}{}{}}
\makeatother
\title{An integral morphism which is a homeomorphism onto its image has property C}
\author{The Stacks Project Authors}
\date{stacks-project:etale-cohomology.tex:local-snapshot-20261011; GNU FDL 1.2 or later}
\begin{document}
\maketitle
\noindent\textit{Editorial note.} All delivered mathematical text is the author's own, copied from the pinned Stacks Project chapter file etale-cohomology.tex; no statement or proof step was rewritten, reordered or removed.\par\medskip
\noindent\textit{Editorial note.} The proof is the author's complete argument, including the reduction to a single point and the assembly of the local pieces into the required V and morphism.\par\medskip
\noindent\textit{Editorial note.} Every labelled reference occurring inside the counted spans is carried in assembly.external\_references as a non-load-bearing navigation pointer (same-chapter ones at their definition site, cross-chapter ones at the invocation inside the counted span), so no load-bearing mathematics is replaced by a stub. No bibliography citation occurs inside the counted spans.\par\medskip
\noindent\textit{Editorial note.} The checked chapter list for same-name collisions is the set of stacks chapters staged in this repository: chow, cohomology, divisors, etale-cohomology, flat, more-algebra, more-morphisms, morphisms, spaces-more-morphisms, varieties together with this card's chapter etale-cohomology.\par\medskip
\noindent\textit{Editorial note.} The preamble is the chapter's own native preamble excerpt already used for this source: the Xy-pic setup (native preamble lines 14-18) and the native theorem declarations with the mathematical macros (native preamble lines 154-217).\par\medskip
\noindent\textit{Allowed standard prerequisites (editorial).}\begin{itemize}
\item Property (C) for a morphism of schemes. Same chapter, the formulation of property (C) delivered verbatim as a context span in this card (source lines 5708-5726): for every etale U over X there is an etale V over Y and a surjective morphism X x\_Y V -> U over X. Supplies the condition verified in the proof.
\item Property (B) for an integral morphism which is a homeomorphism onto its image. Same chapter, the companion result showing that such a morphism has property (B). Supplies the etale covering of the relevant open of Y with which the construction is completed.
\end{itemize}\medskip
\ifdefined\hypertarget\hypertarget{sourcepointer-1}{}\fi
\noindent\textit{Source reference.} Original lemma 04DR: sources/stacks-etale-cohomology/pinned-original-source.tex, source lines 6026--6026. Navigation pointer invoked at proof line 6084; the counted span contains the pointer only and the cited result is not reproduced inside it, so the pointer is registered as a non-load-bearing navigation stub in the established form for Stacks material. No bibliography citation is involved.\par\medskip
\noindent
Let $f : X \to Y$ be a morphism of schemes.
Here is a list of conditions we will consider in the following:
\begin{enumerate}
\item[(A)] For every \'etale morphism $U \to X$ and $u \in U$ there exist
an \'etale morphism $V \to Y$ and a disjoint union decomposition
$X \times_Y V = W \amalg W'$ and a morphism $h : W \to U$ over $X$
with $u$ in the image of $h$.
\item[(B)] For every $V \to Y$ \'etale, and every \'etale covering
$\{U_i \to X \times_Y V\}$ there exists an \'etale covering
$\{V_j \to V\}$ such that for each $j$ we have
$X \times_Y V_j = \coprod W_{ij}$ where $W_{ij} \to X \times_Y V$
factors through $U_i \to X \times_Y V$ for some $i$.
\item[(C)] For every $U \to X$ \'etale, there exists a $V \to Y$ \'etale
and a surjective morphism $X \times_Y V \to U$ over $X$.
\end{enumerate}
It turns out that each of these properties has meaning in terms of
the behaviour of the functor $f_{small, *}$. We will work this
out in the next few sections.

\begin{lemma}
\label{lemma-integral-homeo-onto-image-C}
Let $f : X \to Y$ be an integral morphism of schemes which defines
a homeomorphism of $X$ with a closed subset of $Y$.
Then property (C) holds.
\end{lemma}

\begin{proof}
Let $g : U \to X$ be an \'etale morphism. We need to find an object
$V \to Y$ of $Y_\etale$ and a surjective morphism $X \times_Y V \to U$
over $X$. Suppose that for every $u \in U$ we can find an object
$V_u \to Y$ of $Y_\etale$ and a morphism $h_u : X \times_Y V_u \to U$
over $X$ with $u \in \Im(h_u)$. Then we can take $V = \coprod V_u$
and $h = \coprod h_u$ and we win. Hence given a point
$u \in U$ we find a pair $(V_u, h_u)$ as above. To do this we may
shrink $U$ and assume that $U$ is affine. In this case
$g : U \to X$ is locally quasi-finite. Let
$g^{-1}(g(\{u\})) = \{u, u_2, \ldots, u_n\}$. Since there are no
specializations $u_i \leadsto u$ we may replace $U$ by an affine neighbourhood
so that $g^{-1}(g(\{u\})) = \{u\}$.

\medskip\noindent
The image $g(U) \subset X$ is open,
hence $f(g(U))$ is locally closed in $Y$. Choose an open $V \subset Y$ such
that $f(g(U)) = f(X) \cap V$. It follows that $g$ factors through
$X \times_Y V$ and that the resulting $\{U \to X \times_Y V\}$ is an \'etale
covering. Since $f$ has property (B) , see
Lemma \ref{lemma-integral-B},
we see that there exists an \'etale covering $\{V_j \to V\}$ such that
$X \times_Y V_j \to X \times_Y V$ factor through $U$.
This implies that $V' = \coprod V_j$ is \'etale over $Y$ and that there is a
morphism $h : X \times_Y V' \to U$ whose image
surjects onto $g(U)$. Since $u$ is the only point in its fibre it must
be in the image of $h$ and we win.
\end{proof}

\end{document}
